\documentclass[11pt,a4paper]{article}

% ------------------------------------------------------------------
% Packages
% ------------------------------------------------------------------
\usepackage[utf8]{inputenc}
\usepackage[T1]{fontenc}
\usepackage{lmodern}
\usepackage[margin=2.5cm]{geometry}
\usepackage{amsmath,amssymb,mathtools}
\usepackage{booktabs}
\usepackage{graphicx}
\usepackage{xcolor}
\usepackage{listings}
\usepackage{caption}
\usepackage{algorithm}
\usepackage{algpseudocode}
\usepackage[colorlinks=true,linkcolor=blue!60!black,citecolor=blue!60!black,urlcolor=blue!60!black]{hyperref}

% ------------------------------------------------------------------
% Macros
% ------------------------------------------------------------------
\newcommand{\E}{\mathbb{E}}
\newcommand{\Q}{\mathbb{Q}}
\newcommand{\dd}{\mathrm{d}}
\newcommand{\code}[1]{\texttt{#1}}
\newcommand{\AI}{\mathrm{AI}}

% Images extracted from Products/callable_bond/callable_bond_pricer.ipynb by make_figures.py ; a framed placeholder
% is printed while a file is missing.
\newcommand{\figfile}[2]{%
  \IfFileExists{figures/#1}{\includegraphics[width=#2]{figures/#1}}%
  {\fbox{\parbox[c][4cm][c]{#2}{\centering\small missing figure:\\ \texttt{figures/#1}}}}}

\lstset{basicstyle=\ttfamily\footnotesize, keywordstyle=\color{blue!70!black}, commentstyle=\color{gray},
        stringstyle=\color{teal}, frame=single, breaklines=true, language=Python, showstringspaces=false}

% ------------------------------------------------------------------
\title{Callable Bonds under the Hull--White Model:\\
European, Bermudan and American Calls with Three Numerical Methods}
\author{Elias Benguedda}
\date{November 2026}

\begin{document}
\emergencystretch=3em
\maketitle

\begin{abstract}
We price fixed-rate callable bonds with issuer default risk under a one-factor Hull--White short rate fitted to the
euro area AAA government curve of 5 October 2026 and calibrated to co-terminal swaptions ($a = 0.0863$,
$\sigma = 111.9$bp, RMSE 0.76bp of normal volatility), with a deterministic default intensity bootstrapped from CDS
spreads. The European, Bermudan and American calls differ only by the set of dates where the issuer's choice
$V \leftarrow \min(V, K + \AI)$ is applied, and three independent methods solve the same problem: a Crank--Nicolson
PDE, a Hull--White trinomial tree with unequal time steps, and Monte Carlo with a Longstaff--Schwartz call policy
estimated and applied on independent paths. On a seasoned 4.75\% bond maturing in 2036 and callable at par from
2029, the three methods agree within 0.2 cent per 100 (tree) and 1.5 standard errors (Monte Carlo); the PDE
reprices the straight bond to $3\times10^{-7}$ and the European callable bond to $10^{-4}$ of Jamshidian's closed
form. The issuer's option is worth 1.20, 1.74 and 1.85 per 100 for the European, Bermudan and American calls (6.74
to 7.40 for the same bond without default), the OAS at a clean price of 98.50 falls from 139bp (Z-spread) to 106,
93 and 90bp, and the effective convexity turns from $+67$ to between $-70$ and $-92$. At identical calibration to
European swaptions, the mean reversion moves the American price by up to 0.44 per 100 and the European by 0.01:
the value of the right to choose the call date is model risk.
\end{abstract}

\tableofcontents
\newpage

% ==================================================================
\section{Introduction}
% ==================================================================
A callable bond is a fixed-rate bond that its issuer may redeem before maturity at a contractual call price plus
accrued interest. Issuers call when rates, or their own credit spread, have fallen enough to refinance more
cheaply; the bondholder is short this option and is compensated by a higher coupon. From the holder's side,
\begin{equation}\label{eq:decomposition}
\text{callable bond} = \text{straight bond} - \text{issuer call option}.
\end{equation}
Three exercise styles are common: \emph{European} (one call date, at the end of the call protection),
\emph{Bermudan} (every coupon date after the protection, the most frequent structure for corporate and agency
bonds) and \emph{American} (any time in the call period, typical of US corporate bonds with make-whole or par
calls). The call is an option on a coupon bond, hence on rates: a short-rate model fitted to the yield curve is the
natural framework \cite{hull1990,brigo2006,andersen2010}, and the one-factor Hull--White model is the market
standard for this product because zero-coupon bonds and their options are in closed form.

This report documents a complete pricer: from a term sheet in calendar dates to prices, spreads, yields, risk
measures and the exercise policy, for the three styles. Its distinctive feature is that every number is obtained
by three independent methods (PDE, tree, Monte Carlo) and checked against closed forms where they exist. Section
\ref{sec:model} sets out the model and the product, Section~\ref{sec:methods} the three numerical methods,
Section~\ref{sec:validation} the validation, Section~\ref{sec:results} the prices, spreads and risk of an
illustrative bond, Section~\ref{sec:exercise} the exercise policy and Section~\ref{sec:modelrisk} the model risk.

% ==================================================================
\section{Model and Product}\label{sec:model}
% ==================================================================

\subsection{Hull--White short rate}
Under the risk-neutral measure $\Q$, the short rate follows \cite{hull1990}
\begin{equation}
\dd r_t = \big(\theta(t) - a r_t\big)\,\dd t + \sigma\,\dd W_t,
\end{equation}
written $r_t = x_t + \alpha(t)$ with the Ornstein--Uhlenbeck state $\dd x_t = -a x_t\,\dd t + \sigma\,\dd W_t$,
$x_0 = 0$, and the deterministic shift
\begin{equation}
\alpha(t) = f(0,t) + \frac{\sigma^2}{2a^2}\big(1 - e^{-at}\big)^2,
\end{equation}
which makes the model reprice every discount factor $P(0,T)$ of the initial curve, $f(0,t)$ being the
instantaneous forward rate \cite{brigo2006}. Zero-coupon bonds are exponential-affine in the state:
\begin{equation}\label{eq:zcb}
\begin{aligned}
P(t,T;x) &= \frac{P(0,T)}{P(0,t)}\exp\!\big(-B(t,T)\,x - \Gamma(t,T)\big), \qquad
B(t,T) = \frac{1 - e^{-a(T-t)}}{a},\\
\Gamma(t,T) &= \frac{\sigma^2}{2a^2}B(t,T)\big(1-e^{-at}\big)^2 + \frac{\sigma^2}{4a}\big(1-e^{-2at}\big)B(t,T)^2.
\end{aligned}
\end{equation}
A higher state lowers every bond price: the issuer calls when $x$, hence $r$, is low.

\subsection{Issuer default}
Default is the first jump of a Cox process with deterministic, piecewise constant intensity $\lambda(t)$,
independent of rates \cite{lando1998,duffie1999}; the survival probability is
$Q(0,t) = \exp(-\int_0^t \lambda)$. At default the holder receives $RN$ (recovery of par). The intensities are
bootstrapped sequentially from CDS par spreads (quarterly premiums with accrual on default), each one repricing its
CDS exactly. With $\lambda \equiv 0$ the bond is default-free (government or agency).

\subsection{The product}
The bond pays coupons $c_i$ at dates $T_1 < \dots < T_n$ and the notional $N$ at $T_n$. The issuer may redeem it on
the call set $\mathcal{T}$ at the clean call price $K(t)$ plus accrued interest $\AI(t)$; on a coupon date the holder
receives the coupon and the call price. The three styles are
\begin{equation}
\mathcal{T}^{\mathrm{Eur}} = \{T_c\}, \qquad
\mathcal{T}^{\mathrm{Ber}} = \{T_i : T_c \le T_i < T_n\}, \qquad
\mathcal{T}^{\mathrm{Am}} = [T_c, T_n],
\end{equation}
with $T_c$ the end of the call protection. The term sheet is given in calendar dates: model time is ACT/365F from
the valuation date, coupon dates roll back from the maturity date, each regular coupon pays $Nc/f$ (ACT/ACT ICMA) and
$\AI(t)$ is the running coupon pro rata of its period. A seasoned bond carries accrued interest at the valuation date;
the model gives the \emph{dirty} price, and the market quotes the \emph{clean} price, dirty minus $\AI(0)$.

\subsection{Pricing problem}
The issuer minimises the value of its debt. Before default, the value $V(t,x)$ of the bond is the value of an
optimal stopping problem, $V(0,0) = \inf_{\tau \in \mathcal{T}} \E\big[\text{discounted flows until } \tau\big]$,
and by Feynman--Kac it solves, between payment and call dates,
\begin{equation}\label{eq:pde}
\partial_t V - a x\,\partial_x V + \tfrac12\sigma^2\,\partial_{xx} V - \big(x + \alpha(t) + \lambda(t)\big)V
+ \lambda(t) R N = 0,
\end{equation}
with $V(T_n^-, x) = N + c_n$ and, at a date $t$,
\begin{equation}\label{eq:jump}
V(t^-, x) = \min\big(V(t^+, x),\ K(t) + \AI(t)\big)\,\mathbf{1}_{t \in \mathcal{T}} + V(t^+,x)\,\mathbf{1}_{t \notin \mathcal{T}} + c(t),
\end{equation}
where $c(t)$ is the coupon paid at $t$, if any. Default enters as an extra discount rate $\lambda$ and a recovery
source term. The three styles differ only by $\mathcal{T}$: one solver serves all of them. On a call date, the
\emph{critical state} $x^*(t) = \max\{x : V(t^+,x) \ge K(t) + \AI(t)\}$ separates calling (below) from continuing,
and $r^*(t) = x^*(t) + \alpha(t)$ is the critical short rate.

% ==================================================================
\section{Numerical Methods}\label{sec:methods}
% ==================================================================
The three methods share the time grid: the valuation date, every coupon, call and CDS pillar date, and between two
of them equal steps of at most $1/100$ year (975 steps for the bond of Section~\ref{sec:results}). On the American
call period, the issuer may call at every node.

\subsection{Crank--Nicolson PDE}
Since $\alpha$ and $\lambda$ do not depend on $x$, on each step $[t_{k-1}, t_k]$ the solution is
\begin{equation}
V(t_{k-1}, x) = e^{-\int_{t_{k-1}}^{t_k}\alpha}\,\frac{Q(0,t_k)}{Q(0,t_{k-1})}\,U(t_{k-1},x)
+ RN\Big(1 - \frac{Q(0,t_k)}{Q(0,t_{k-1})}\Big) P\big(t_{k-1}, t_{k-1} + \tfrac{\Delta t}{2}; x\big),
\end{equation}
where $U$ solves the time-homogeneous equation $U_t - a x U_x + \frac12\sigma^2 U_{xx} - xU = 0$ with
$U(t_k) = V(t_k)$, and the recovery is paid at mid-step. The deterministic discounting $\int\alpha$ is exact (closed
form), and the operator of $U$ is constant, so each step is one tridiagonal solve of the $\theta$-scheme
\begin{equation}
(I - \theta\,\Delta t\,L)\,U^{k-1} = \big(I + (1-\theta)\,\Delta t\,L\big)\,U^{k},
\end{equation}
with $L$ the central-difference discretisation of $-ax\partial_x + \frac12\sigma^2\partial_{xx} - x$, one-sided and
upwind at the boundaries, where the drift points inward. The grid is uniform in $x$ over $\pm 6$ standard deviations
of $x(T_n)$ with 401 nodes. $\theta = \frac12$ (Crank--Nicolson) except for two fully implicit steps after each
discrete call date (Rannacher smoothing \cite{rannacher1984}), which damp the oscillations created by the kink of
the $\min$ in~\eqref{eq:jump}. The PDE returns $V(0,x)$ for every initial state, hence the price as a function of
$r(0)$, and the critical states $x^*(t)$.

\subsection{Trinomial tree}
The Hull--White tree \cite{hull1994} is built on the state $x$ and on the same unequal time grid. Level $k+1$ has its
own spacing, set by the variance of the step that leads to it,
\begin{equation}
\Delta x_{k+1} = \sqrt{3V_k}, \qquad V_k = \sigma^2\,\frac{1 - e^{-2ah_k}}{2a}, \qquad h_k = t_{k+1} - t_k.
\end{equation}
A node $x$ of level $k$ branches to the node $j\,\Delta x_{k+1}$ nearest to its conditional mean $x e^{-ah_k}$ and
to its two neighbours. With $\eta = (xe^{-ah_k} - j\Delta x_{k+1})/\Delta x_{k+1} \in [-\frac12, \frac12]$ and
$m = \frac13 + \eta^2$, the probabilities matching the mean and variance are
\begin{equation}
p_u = \frac{m + \eta}{2}, \qquad p_m = 1 - m, \qquad p_d = \frac{m - \eta}{2},
\end{equation}
all positive ($p_m \ge \frac{5}{12}$, $p_u, p_d \ge \frac{1}{24}$) whatever the step. Levels are truncated at
$\pm 7$ standard deviations of $x(t)$ (179 nodes at maturity). The shift is not taken from the closed form: the
discount factor of a node over a step is $e^{-(\alpha_k + x)h_k}$ and $\alpha_k$ is solved by forward induction of
the Arrow--Debreu prices $q_{k}(x)$,
\begin{equation}
\sum_{x} q_k(x)\,e^{-(\alpha_k + x)h_k} = P(0, t_{k+1}), \qquad
q_{k+1}(x') = \sum_{x \to x'} q_k(x)\,e^{-(\alpha_k + x)h_k}\,p_{x \to x'},
\end{equation}
so that the tree reprices every discount factor of the curve exactly. Default and recovery enter as in the PDE. The
backward induction applies~\eqref{eq:jump} at the nodes.

\subsection{Monte Carlo and Longstaff--Schwartz}
The pair $(x_t, \int_0^t x_s\,\dd s)$ is Gaussian and simulated exactly on the grid of coupon, call and (with
recovery) monthly dates, which gives the state and the bank-account discount factor $D(0,t) = e^{-\int_0^t r}$ on
each path without discretisation bias. The call policy is that of Longstaff and Schwartz \cite{longstaff2001},
adapted to an option held by the issuer (Algorithm~\ref{alg:lsm}). Two choices matter:
\begin{itemize}
\item \emph{Regressors.} Besides $1, z, z^2, z^3$ ($z = x/\mathrm{sd}(x_t)$), the closed-form value at $t$ of the
remaining straight flows, coupons, redemption and recovery, $S_t(x)$, computed with~\eqref{eq:zcb}. Since the call
can only lower the bond, the continuation value is below $S_t(x)$: only the paths with $S_t(x) > K + \AI$ can be
called, and only they are regressed.
\item \emph{Out-of-sample policy.} The coefficients are estimated on 50{,}000 paths and applied to 50{,}000
independent ones. The true price is the infimum over call policies; any feasible policy gives a higher bond price, so
the out-of-sample price is an \emph{upper bound} of the true price up to Monte Carlo noise, whereas the in-sample
price reuses the paths that defined the policy.
\end{itemize}
The American call is approximated by monthly call dates. Antithetic variates reduce the variance; the standard
error is computed on the pair averages.

\begin{algorithm}[h]
\caption{Longstaff--Schwartz for a callable bond (issuer's call)}
\label{alg:lsm}
\begin{algorithmic}[1]
\State simulate $(x, D)$ on the dates $t_1 < \dots < t_M$; \quad $\Pi \gets 0$ \Comment{PV at 0 of the flows after $t_k$}
\For{$k = M, \dots, 1$}
  \If{$t_k$ is a call date, strike $\kappa_k = K(t_k) + \AI(t_k)$}
    \State candidates: paths with $S_{t_k}(x) > \kappa_k$
    \State fit $\beta_k$: regress $\Pi / (D(0,t_k)Q(0,t_k))$ on $\varphi(x) = (1,z,z^2,z^3,S_{t_k}(x))$ over the candidates
    \State call on the candidates where $\varphi(x)^\top\beta_k > \kappa_k$: \quad $\Pi \gets \kappa_k\,D(0,t_k)\,Q(0,t_k)$
  \EndIf
  \State $\Pi \gets \Pi + c(t_k)\,D(0,t_k)\,Q(0,t_k) + RN\,\big(Q(0,t_{k-1}) - Q(0,t_k)\big)\tfrac12\big(D(0,t_{k-1}) + D(0,t_k)\big)$
\EndFor
\State re-simulate independent paths and repeat lines 2--8 with the fitted $\beta_k$; \quad price $= \overline{\Pi}$
\end{algorithmic}
\end{algorithm}

\subsection{Closed forms used as benchmarks}
\emph{Straight bond.} With default independent of rates,
\begin{equation}
B_0 = \sum_i c_i P(0,T_i)Q(0,T_i) + RN\int_0^{T_n} P(0,t)\,\big(-\dd Q(0,t)\big).
\end{equation}

\emph{European call, zero recovery} \cite{jamshidian1989}. At $T_c$, given survival, the bond is worth
$\sum_{T_i > T_c} \tilde c_i P(T_c, T_i; x)$ with $\tilde c_i = c_i Q(0,T_i)/Q(0,T_c)$, decreasing in $x$. With $x^*$
such that this sum equals the strike $\kappa = K(T_c) + \AI(T_c)$ and $\kappa_i = P(T_c,T_i;x^*)$, the option is a
portfolio of options on zero-coupon bonds, held only if the issuer survives to $T_c$:
\begin{equation}
V^{\mathrm{Eur}}_0 = B_0 - Q(0,T_c)\sum_{T_i > T_c}\tilde c_i\,\mathrm{ZBC}(0, T_c, T_i, \kappa_i),
\end{equation}
where $\mathrm{ZBC}$ is the Hull--White zero-coupon bond call
\begin{equation}
\begin{aligned}
\mathrm{ZBC}(0,T_c,T_i,\kappa_i) &= P(0,T_i)\Phi(h) - \kappa_i P(0,T_c)\Phi(h - \sigma_p),\\
\sigma_p &= \sigma B(T_c,T_i)\sqrt{\tfrac{1-e^{-2aT_c}}{2a}}, \qquad
h = \frac{1}{\sigma_p}\ln\frac{P(0,T_i)}{\kappa_i P(0,T_c)} + \frac{\sigma_p}{2}.
\end{aligned}
\end{equation}

\subsection{Spreads, yields and risk}
\begin{itemize}
\item \emph{OAS} \cite{kalotay1993}: the constant spread $s$ such that discounting at $r_t + s$ (no explicit default)
reprices the dirty market price with the call; the \emph{Z-spread} is the same spread without the call. The
difference measures the share of the yield that pays for the short option.
\item \emph{Yield to worst}: the lowest of the yields to each redemption date (call dates at $K + \AI$, and maturity
at $N$), compounded at the coupon frequency.
\item \emph{Effective duration, convexity, DV01}: full repricing under $\pm 10$bp parallel shifts of the zero curve,
the Hull--White shift refitted to each curve, $a$ and $\sigma$ fixed. \emph{Vega}: per bp of $\sigma$.
\emph{CS01}: per bp of parallel shift of the hazard rates.
\item \emph{Key-rate durations} \cite{ho1992}: triangular shifts of each pillar zero rate, linear to the neighbouring
pillars, applied to the base curve sampled weekly and interpolated linearly. Bumping a single pillar of the
monotone convex curve instead does not work: that interpolation is not linear in its inputs, and the single-pillar
bumps add up to 7.11 instead of 7.56 for the straight bond. The triangular shifts add up to a parallel shift exactly.
\end{itemize}

% ==================================================================
\section{Market Data and Calibration}\label{sec:data}
% ==================================================================
\emph{Curve.} The risk-free curve is the ECB euro area AAA government spot curve of 5 October 2026 (18 pillars from
3 months to 30 years; zero rates 2.85\% at 1 year, 3.06\% at 3 years, 3.50\% at 10 years), interpolated by monotone
convex \cite{benguedda2026yield}.

\emph{Volatility.} The call is an option on the co-terminal swaps, ending at the bond maturity. $(a, \sigma)$ are
calibrated to the ATM co-terminal European receiver swaptions $T \times (9.70 - T)$, $T = 1,\dots,5$, minimising the
pricing errors divided by the market vega, i.e. errors in normal volatility. The normal volatilities are an
\emph{illustrative} EUR grid (no free quotes): about 77 to 79bp. Result: $a = 0.0863$, $\sigma = 111.9$bp, RMSE
0.76bp, maximum error 1.17bp.

\emph{Credit.} The CDS spreads of the issuer are \emph{illustrative} for a BBB name: 60, 85, 110, 125 and 140bp at
1, 3, 5, 7 and 10 years, $R = 40\%$. The bootstrapped hazard rates are 1.00\%, 1.63\%, 2.53\%, 2.84\% and 3.12\%; the
10-year survival probability is 78.4\%.

\emph{Term sheet} (illustrative). Issued 15/06/2024, maturity 15/06/2036, 4.75\% annual coupon, notional 100, non-call
until 15/06/2029, callable at par. At the valuation date the maturity is 9.70 years, the first call 2.70 years and the
accrued interest 1.4575. The European bond is callable on 15/06/2029 only, the Bermudan on the seven coupon dates
2029 to 2035, the American at any time from 15/06/2029 to maturity.

% ==================================================================
\section{Validation}\label{sec:validation}
% ==================================================================

\subsection{Closed forms}
Table~\ref{tab:closed} compares the three methods with the closed forms. The PDE and the tree reprice the straight
bond to machine precision (the tree by construction of $\alpha_k$). The European call is recovered to $10^{-4}$ by the
PDE and to $1.4\times10^{-3}$ by the tree, and Monte Carlo is within 1.7 standard errors in every case.

\begin{table}[h]
\centering
\begin{tabular}{@{}llrrrrr@{}}
\toprule
& \textbf{Case} & \textbf{Closed form} & \textbf{PDE} & \textbf{Tree} & \textbf{LSM} & \textbf{LSM s.e.} \\
\midrule
Straight & default-free & 111.5565 & $3\times10^{-7}$ & $-1\times10^{-12}$ & 0.0059 & 0.0078 \\
         & credit, $R = 0$ & 92.9239 & $3\times10^{-7}$ & $-1\times10^{-12}$ & 0.0047 & 0.0062 \\
         & credit, $R = 40\%$ & 99.9402 & $2\times10^{-7}$ & $-2\times10^{-7}$ & 0.0110 & 0.0065 \\
European & default-free & 104.8151 & $-2\times10^{-5}$ & $-0.0014$ & $-0.0066$ & 0.0174 \\
callable & credit, $R = 0$ & 92.6967 & $-9\times10^{-5}$ & $-0.0009$ & 0.0090 & 0.0077 \\
\bottomrule
\end{tabular}
\caption{Method minus closed form (dirty prices, per 100). Jamshidian's formula only holds with zero recovery.}
\label{tab:closed}
\end{table}

\subsection{Agreement of the three methods}
Table~\ref{tab:methods} prices the four bonds with the three methods. The tree is within 0.16 cent of the PDE;
Monte Carlo is 1.0 to 1.5 standard errors below it (the four estimates use the same seed, hence share their noise).
For the American bond, Monte Carlo prices monthly call dates, slightly above the continuous American
(Table~\ref{tab:bermam}). Times are 0.3 to 0.55 second for the tree and the PDE; Monte Carlo takes
6.3 seconds for the American (about 90 call dates, each with a closed-form regressor on every path).

\begin{table}[h]
\centering
\begin{tabular}{@{}lrrrrrrr@{}}
\toprule
\textbf{Bond} & \textbf{PDE} & \textbf{Tree} & \textbf{LSM} & \textbf{s.e.} & \textbf{$\frac{\text{LSM}-\text{PDE}}{\text{s.e.}}$} & \textbf{PDE (s)} & \textbf{LSM (s)} \\
\midrule
Straight & 99.9402 & 99.9402 & 99.9512 & 0.0065 & $+1.68$ & 0.46 & 0.41 \\
European & 98.7382 & 98.7388 & 98.7149 & 0.0158 & $-1.47$ & 0.46 & 0.49 \\
Bermudan & 98.2050 & 98.2034 & 98.1854 & 0.0135 & $-1.46$ & 0.46 & 0.82 \\
American & 98.0890 & 98.0877 & 98.0758 & 0.0132 & $-0.99$ & 0.54 & 6.29 \\
\bottomrule
\end{tabular}
\caption{Dirty prices per 100 by method (credit, $R = 40\%$). LSM: 50{,}000 out-of-sample paths, monthly calls for the American.}
\label{tab:methods}
\end{table}

\subsection{Longstaff--Schwartz bias}
Over five independent seeds, the Bermudan price is 98.1985 out-of-sample and 98.2121 in-sample, against 98.2050 for
the PDE: $-0.0065$ and $+0.0071$, both within the standard error of a five-seed mean (about 0.006). The regression
policy is therefore close to optimal; the out-of-sample price is the one to report, since it values a policy the
issuer could follow.

\subsection{From Bermudan to American}
A bond callable every $1/m$ year from the first call date must converge to the American bond as $m\to\infty$, and be
worth less with more call dates. Table~\ref{tab:bermam} confirms both, with the tree and the PDE agreeing at each
frequency.

\begin{table}[h]
\centering
\begin{tabular}{@{}lrrrrrr@{}}
\toprule
\textbf{Call dates} & annual & semi-annual & quarterly & monthly & weekly & American \\
\midrule
PDE  & 98.2050 & 98.1491 & 98.1196 & 98.0988 & 98.0918 & 98.0890 \\
Tree & 98.2034 & 98.1480 & 98.1178 & 98.0970 & 98.0892 & 98.0877 \\
\bottomrule
\end{tabular}
\caption{Bermudan bonds with denser call dates converge to the American bond (dirty, per 100).}
\label{tab:bermam}
\end{table}

\subsection{Grid convergence}
Figure~\ref{fig:convergence} measures the error against a fine PDE (3{,}201 nodes, 800 steps per year), the space
grid of the PDE being refined with the time grid. Crank--Nicolson with Rannacher steps converges fast for the European
and Bermudan bonds (below $10^{-4}$ from 50 steps per year). The American converges at first order in time, as any
scheme where the $\min$ is applied at every node. The tree error is not monotone (the position of the nodes relative
to the strike changes with the grid), but stays below 0.2 cent from 100 steps per year. The default setting, 100
steps per year, gives errors below 0.12 cent per 100 for the three styles.

\begin{figure}[h]
\centering
\figfile{convergence.png}{\textwidth}
\caption{Absolute error against a fine PDE as a function of the time steps per year.}
\label{fig:convergence}
\end{figure}

% ==================================================================
\section{Prices, Spreads and Risk}\label{sec:results}
% ==================================================================

\subsection{Prices}
Table~\ref{tab:prices} gives the model prices. The more call dates, the more valuable the issuer's option: it is worth
1.20 per 100 for the European call, 1.74 for the Bermudan and 1.85 for the American. The first call date carries most
of the value; the Bermudan adds the right to wait, and the American only adds the dates between coupons.

Default lowers the option value by a factor of four to six: without default, the same bonds are worth 104.82, 104.29 and
104.16 against 111.56 for the straight bond, an option of 6.74 to 7.40. With credit risk, the bond trades below par,
further from the call price, and the call is out of the money; moreover the issuer can only call if it survives.

\begin{table}[h]
\centering
\begin{tabular}{@{}lrrrr@{}}
\toprule
& \textbf{Straight} & \textbf{European} & \textbf{Bermudan} & \textbf{American} \\
\midrule
Dirty price & 99.9402 & 98.7382 & 98.2050 & 98.0890 \\
Clean price & 98.4827 & 97.2807 & 96.7475 & 96.6314 \\
Issuer option & -- & 1.2020 & 1.7352 & 1.8513 \\
Default-free dirty price & 111.5565 & 104.8151 & 104.2878 & 104.1586 \\
Default-free issuer option & -- & 6.7415 & 7.2687 & 7.3979 \\
\bottomrule
\end{tabular}
\caption{Model prices (PDE, per 100; accrued interest 1.4575).}
\label{tab:prices}
\end{table}

\subsection{Spreads and yields from a quoted price}
In practice the clean price is observed and the spreads are solved for. Table~\ref{tab:market} uses an illustrative
quote of 98.50. The Z-spread, 139bp, ignores the call; the OAS strips it out and is lower for every style, the more so
with more call dates: 106bp (European), 93bp (Bermudan) and 90bp (American). Part of the extra yield of a callable bond
pays for the option the holder has sold, not for credit.

Below par, the worst case for the holder is the bond running to maturity, and the yield to worst equals the yield to
maturity (4.94\%). Above par it is an early call: at a clean price of 104 the yield to maturity is 4.23\% but the yield
to the first call is 3.17\% (Figure~\ref{fig:yields}).

\begin{table}[h]
\centering
\begin{tabular}{@{}lrrr@{}}
\toprule
& \textbf{European} & \textbf{Bermudan} & \textbf{American} \\
\midrule
Z-spread (bp) & 138.7 & 138.7 & 138.7 \\
OAS (bp) & 106.3 & 92.8 & 89.9 \\
Yield to maturity (\%) & 4.941 & 4.941 & 4.941 \\
Yield to worst (\%) & 4.941 & 4.941 & 4.941 \\
\bottomrule
\end{tabular}
\caption{Spreads and yields at a quoted clean price of 98.50 (dirty 99.96).}
\label{tab:market}
\end{table}

\begin{figure}[h]
\centering
\figfile{yields.png}{0.8\textwidth}
\caption{Yield to each redemption date of the Bermudan bond at three clean prices; circled: the yield to worst.}
\label{fig:yields}
\end{figure}

\subsection{Interest rate, volatility and credit risk}
Table~\ref{tab:risk} gives the risk measures. The callable bonds have a shorter duration than the straight bond (5.9
to 6.0 against 7.6) and a \emph{negative} convexity: when rates fall, the call gets closer, the expected life
shortens and the price is capped near the call price (Figure~\ref{fig:profile}). The vega is negative, since the
holder is short the option, and grows with the number of call dates. The CS01 is smaller than that of the straight
bond: a wider spread pushes the call further out of the money, which offsets part of the loss.

The three methods agree on duration, DV01, vega and CS01 to the second or third digit (e.g. Bermudan duration 5.927,
5.926 and 5.920 for the PDE, the tree and Monte Carlo, vega $-0.0271$, $-0.0272$, $-0.0270$). The effective
convexity, a second difference, depends on the exercise discretisation: the tree ($-137$ for the European) and Monte
Carlo ($-110$ for the Bermudan, policy re-estimated in each scenario) are noisy; the PDE values are the ones to use.

\begin{table}[h]
\centering
\begin{tabular}{@{}lrrrr@{}}
\toprule
& \textbf{Straight} & \textbf{European} & \textbf{Bermudan} & \textbf{American} \\
\midrule
Effective duration & 7.560 & 6.044 & 5.927 & 5.894 \\
Effective convexity & 66.5 & $-92.0$ & $-69.9$ & $-70.6$ \\
DV01 (per 100) & 0.0756 & 0.0597 & 0.0582 & 0.0578 \\
Vega (per bp of $\sigma$) & 0.0000 & $-0.0199$ & $-0.0271$ & $-0.0285$ \\
CS01 (per bp of hazard) & $-0.0451$ & $-0.0352$ & $-0.0342$ & $-0.0339$ \\
\bottomrule
\end{tabular}
\caption{Effective risk measures (PDE, $\pm 10$bp shifts, Hull--White refitted).}
\label{tab:risk}
\end{table}

\begin{figure}[h]
\centering
\figfile{price_profile.png}{\textwidth}
\caption{Left: $V(0,x)$ as a function of the initial short rate, from one PDE solve (today: dotted). Right: full
repricing under parallel shifts of the curve. The callable bonds are capped near the call price when rates fall.}
\label{fig:profile}
\end{figure}

The key-rate durations (Figure~\ref{fig:krd}) add up to the effective duration (7.5595, 6.0429, 5.9205 and 5.8996
against 7.5595, 6.0429, 5.9200 and 5.8954 for 1bp parallel shifts). The straight bond's risk sits on the 9- and
10-year pillars, near its maturity. The call moves part of it to the 2- to 4-year pillars around the first call date:
a callable bond behaves partly like a bond maturing at the call. The hedge of a callable bond therefore involves
short-dated instruments that a yield-to-maturity duration ignores.

\begin{figure}[h]
\centering
\figfile{key_rate_durations.png}{0.85\textwidth}
\caption{Key-rate durations (triangular shifts, PDE).}
\label{fig:krd}
\end{figure}

% ==================================================================
\section{Exercise Policy}\label{sec:exercise}
% ==================================================================

\subsection{Exercise boundary}
The issuer calls when the short rate is below the critical rate $r^*(t)$ (Figure~\ref{fig:boundary}). On the first
call date, $r^*$ is 2.20\% for the European bond, 1.54\% for the Bermudan and 1.03\% for the American. With more
future call dates the issuer gains by waiting, and only calls for a deeper fall in rates; the European issuer, with
a single chance, calls at the highest threshold. The boundary rises towards maturity (1.28\% after 6.2 years, 2.44\%
about a month before maturity for the American; 2.55\% on the last Bermudan date): with fewer coupons left to save, a
smaller rate drop is enough. All thresholds are well below the forward curve (3.2\% to 4.1\%): the call is out of
the money at today's rates, because the bond, after its credit spread, is worth less than par. The tree gives the
same Bermudan thresholds within one node.

\begin{figure}[h]
\centering
\figfile{boundary.png}{0.85\textwidth}
\caption{Critical short rate $r^*(t)$ below which the issuer calls (PDE; crosses: tree).}
\label{fig:boundary}
\end{figure}

\subsection{Call probabilities and expected life}
The Monte Carlo paths give the probability of a call on each date given survival (Figure~\ref{fig:callprob},
Table~\ref{tab:life}). The European bond is called with probability 28\%. The Bermudan bond is called with
probability 43\% in total, 16\% on the first call date and less than 7\% on each later one. The American issuer calls
more often in total (46\%) but spreads the calls over the call period, waiting for lower rates, so the expected life
(7.60 years) is not shorter than for the Bermudan (7.59 years).

\begin{figure}[h]
\centering
\figfile{call_probability.png}{\textwidth}
\caption{Left: probability of a call on each date, given survival (European and Bermudan). Right: share of the bonds
still outstanding, all styles (American: monthly call dates).}
\label{fig:callprob}
\end{figure}

\begin{table}[h]
\centering
\begin{tabular}{@{}lrrr@{}}
\toprule
& \textbf{European} & \textbf{Bermudan} & \textbf{American (monthly)} \\
\midrule
Probability of a call (\%) & 27.6 & 42.5 & 45.6 \\
Expected life (years) & 7.77 & 7.59 & 7.60 \\
\bottomrule
\end{tabular}
\caption{Call statistics given survival (Monte Carlo, 50{,}000 paths); maturity 9.70 years.}
\label{tab:life}
\end{table}

% ==================================================================
\section{Model Risk and Term Sheet Variants}\label{sec:modelrisk}
% ==================================================================

\subsection{Mean reversion}
The co-terminal European swaptions fix one combination of $(a, \sigma)$ for each expiry, not both parameters
independently. Table~\ref{tab:meanrev} fixes $a$, recalibrates $\sigma$ to the same five swaptions and reprices the
bonds. The fit stays acceptable (RMSE below 1bp of normal volatility up to $a = 0.15$). The European bond, an option
close to the calibration set, moves by at most 1.4 cent. The Bermudan and American bonds move by up to 0.32 and 0.44
per 100: the mean reversion sets the correlation between rates at different call dates, hence the value of choosing
the call date, which European swaptions do not see. A higher $a$ decorrelates the call dates, makes the right to
choose worth more and the bond worth less. This spread is model risk at identical calibration; reducing it requires
Bermudan swaption quotes or a second factor.

\begin{table}[h]
\centering
\begin{tabular}{@{}rrrrrr@{}}
\toprule
$a$ & $\sigma$ (bp) & \textbf{RMSE (bp)} & \textbf{European} & \textbf{Bermudan} & \textbf{American} \\
\midrule
0.005 & 77.9 & 0.93 & $+0.014$ & $+0.142$ & $+0.182$ \\
0.020 & 83.5 & 0.88 & $+0.011$ & $+0.117$ & $+0.151$ \\
0.050 & 95.6 & 0.80 & $+0.005$ & $+0.065$ & $+0.085$ \\
\textbf{0.0863} & \textbf{111.9} & \textbf{0.76} & \textbf{98.7382} & \textbf{98.2050} & \textbf{98.0890} \\
0.150 & 144.8 & 0.91 & $-0.006$ & $-0.121$ & $-0.162$ \\
0.250 & 208.1 & 1.58 & $-0.007$ & $-0.324$ & $-0.445$ \\
\bottomrule
\end{tabular}
\caption{Prices with $\sigma$ recalibrated for each mean reversion, as differences to the calibrated model (bold row).}
\label{tab:meanrev}
\end{table}

\subsection{Volatility and credit}
The issuer's option increases with the volatility (Figure~\ref{fig:volcredit}): from 0.01 to 0.03 per 100 at
$\sigma = 30$bp to 3.07 (European), 4.21 (Bermudan) and 4.45 (American) at 200bp. It decreases with the credit
spread: a parallel shift of the hazard rates of $-100$bp raises it to 2.54 to 3.27, $+300$bp lowers it to 0.05 to
0.22, as the bond moves away from the call price. With a deterministic intensity, the model ignores the other reason
to call, a tightening of the issuer's own spread (Section~\ref{sec:limitations}).

\begin{figure}[h]
\centering
\figfile{vol_credit.png}{\textwidth}
\caption{Issuer option value as a function of the Hull--White volatility (left) and of a parallel shift of the hazard
rates (right).}
\label{fig:volcredit}
\end{figure}

\subsection{Term sheet variants}
Table~\ref{tab:variants} prices three common variants of the bond. A step-down call price (102 from 2029, 101 from
2031, par from 2033) makes early calls costly and lowers the option by 0.42 for the European bond, whose only call is
at 102, and by about 0.40 for the other two. Semi-annual coupons (same annual rate) raise the bond, since coupons are
received earlier, and give the Bermudan issuer twice as many call dates. Quarterly calls on the annual-coupon bond
(call price plus accrued interest between coupons) lie between the annual Bermudan and the American. The tree agrees
with the PDE within 0.2 cent on every variant.

\begin{table}[h]
\centering
\begin{tabular}{@{}llrrrr@{}}
\toprule
\textbf{Variant} & \textbf{Style} & \textbf{Dirty} & \textbf{Clean} & \textbf{Issuer option} & \textbf{Tree$-$PDE} \\
\midrule
Par call (base) & European & 98.7382 & 97.2807 & 1.2020 & $+0.0006$ \\
                & Bermudan & 98.2050 & 96.7475 & 1.7352 & $-0.0016$ \\
                & American & 98.0890 & 96.6314 & 1.8513 & $-0.0012$ \\
Step-down 102/101/100 & European & 99.1624 & 97.7048 & 0.7779 & $-0.0014$ \\
                & Bermudan & 98.5970 & 97.1394 & 1.3433 & $-0.0012$ \\
                & American & 98.5088 & 97.0513 & 1.4314 & $-0.0012$ \\
Semi-annual coupons & European & 99.1783 & 97.7247 & 1.2944 & $-0.0020$ \\
                & Bermudan & 98.5841 & 97.1306 & 1.8886 & $-0.0013$ \\
                & American & 98.5235 & 97.0699 & 1.9493 & $-0.0018$ \\
Quarterly calls & Bermudan & 98.1196 & 96.6620 & 1.8207 & $-0.0017$ \\
\bottomrule
\end{tabular}
\caption{Term sheet variants (PDE, per 100).}
\label{tab:variants}
\end{table}

% ==================================================================
\section{Implementation}\label{sec:implementation}
% ==================================================================
\begin{itemize}
\item \code{Products/callable\_bond/callable\_bond\_pricer.py}:\\ \code{CreditCurve} (CDS bootstrap),
\code{CallableBond} and \code{DatedCallableBond} (term sheet), \code{CallableBondPDE}, \code{CallableBondTree},
\code{CallableBondLSM} (one interface: \code{price(bond, callable)}, \code{with\_model}, \code{with\_credit}), the
closed forms, \code{option\_adjusted\_spread}, \code{effective\_risk}, \code{yield\_table},
\code{key\_rate\_durations} and \code{pricing\_report};
\item \code{Fixed\_Income\_Models}: \code{YieldCurve}, \code{HullWhite} (zero-coupon bonds and options, exact
simulation, calibration), \code{RatePricer};
\item \code{Products/callable\_bond/callable\_bond\_pricer.ipynb}: the notebook of this report (about 4 minutes);
\code{make\_figures.py} extracts its figures.
\end{itemize}

\begin{lstlisting}
from callable_bond_pricer import DatedCallableBond, CallableBondPDE, CallableBondTree, CallableBondLSM, pricing_report

bond = DatedCallableBond("2026-10-05", "2036-06-15", 0.0475, first_call_date="2029-06-15",
                         issue_date="2024-06-15", exercise="bermudan")
pde = CallableBondPDE(hw, credit, recovery=0.40)          # hw : calibrated HullWhite, credit : CreditCurve
print(pde.price(bond), CallableBondTree(hw, credit, 0.40).price(bond), CallableBondLSM(hw, credit, 0.40).price(bond))
print(pricing_report(pde, bond.with_exercise("american"), market_price=98.50))   # OAS, YTW, duration, vega, CS01
\end{lstlisting}

% ==================================================================
\section{Limitations and Next Steps}\label{sec:limitations}
% ==================================================================
\begin{itemize}
\item \emph{One factor.} Rates of all maturities are perfectly correlated; the Bermudan and American prices depend on
the mean reversion at identical European calibration (Table~\ref{tab:meanrev}). A two-factor Gaussian model (G2++)
calibrated to Bermudan swaptions would reduce this model risk; the Monte Carlo pricer only needs the model's
simulation and zero-coupon bond prices to extend to it.
\item \emph{Credit.} The intensity is deterministic and independent of rates, so the issuer never calls because its
own spread tightens, and the option is undervalued for issuers with volatile spreads. A stochastic intensity
correlated with rates (CIR or a second Gaussian factor) would capture it.
\item \emph{Market data.} Swaption volatilities and CDS spreads are illustrative; the AAA government curve is used as
risk-free curve, whereas an €STR curve is the standard discounting curve for a dealer.
\item \emph{Conventions.} No business-day adjustment, settlement lag or call notice period; the American call is
monthly in Monte Carlo. A notice period changes the payoff to $\big(K + \AI\big)P(t_n, t_c; x)$ at the notice date
and is a direct extension of the three methods.
\item \emph{Product scope.} Puttable bonds (holder's option, $\max$ instead of $\min$), make-whole calls and
amortising bonds are not covered.
\end{itemize}

% ==================================================================
\section{Conclusion}
% ==================================================================
A single pricing equation, with the issuer's $\min$ applied on the call set, prices European, Bermudan and American
callable bonds under Hull--White with default risk. Three independent methods, a Crank--Nicolson PDE, a trinomial tree
with unequal steps and an out-of-sample Longstaff--Schwartz Monte Carlo, agree within a fraction of a cent and
reprice the closed forms. On an illustrative BBB bond, the issuer's option is worth 1.2 to 1.9 per 100, four to six
times less than for the same default-free bond; it lowers the OAS by 32 to 49bp below the Z-spread, turns the convexity
negative and moves the rate risk towards the first call date. The European price is pinned by the swaption
calibration, but the value of choosing the call date, which separates the Bermudan and American bonds from it,
depends on the mean reversion: a model risk of up to 0.44 per 100 that European swaptions alone cannot remove.

% ==================================================================
\appendix
\section{Notation}
% ==================================================================
\begin{center}
\begin{tabular}{@{}ll@{}}
\toprule
Symbol & Meaning \\
\midrule
$r_t = x_t + \alpha(t)$ & Short rate, Ornstein--Uhlenbeck state, deterministic shift \\
$a$, $\sigma$ & Mean reversion, normal volatility of the short rate \\
$P(t,T;x)$, $B(t,T)$ & Zero-coupon bond price given the state, its sensitivity to $x$ \\
$\lambda(t)$, $Q(0,t)$, $R$ & Default intensity, survival probability, recovery rate \\
$N$, $c_i$, $T_i$, $f$ & Notional, coupons, payment dates, coupon frequency \\
$K(t)$, $\AI(t)$, $\mathcal{T}$, $T_c$ & Clean call price, accrued interest, call set, first call date \\
$V(t,x)$ & Value of the callable bond before default \\
$x^*(t)$, $r^*(t)$ & Critical state and short rate below which the issuer calls \\
$D(0,t)$ & Bank-account discount factor $e^{-\int_0^t r}$ \\
\bottomrule
\end{tabular}
\end{center}

% ==================================================================
\begin{thebibliography}{99}
\bibitem{benguedda2026yield} E. Benguedda, \emph{Yield Curve Construction: Interpolation, Bootstrapping and Curve
Risk}, 2026.
\bibitem{andersen2010} L. Andersen and V. Piterbarg, \emph{Interest Rate Modeling}, Atlantic Financial Press, 2010.
\bibitem{brigo2006} D. Brigo and F. Mercurio, \emph{Interest Rate Models -- Theory and Practice}, 2nd edition,
Springer, 2006.
\bibitem{duffie1999} D. Duffie and K. Singleton, Modeling term structures of defaultable bonds, \emph{Review of
Financial Studies}, 12(4):687--720, 1999.
\bibitem{ho1992} T. S. Y. Ho, Key rate durations: measures of interest rate risks, \emph{Journal of Fixed Income},
2(2):29--44, 1992.
\bibitem{hull1990} J. Hull and A. White, Pricing interest-rate-derivative securities, \emph{Review of Financial
Studies}, 3(4):573--592, 1990.
\bibitem{hull1994} J. Hull and A. White, Numerical procedures for implementing term structure models I: single-factor
models, \emph{Journal of Derivatives}, 2(1):7--16, 1994.
\bibitem{jamshidian1989} F. Jamshidian, An exact bond option formula, \emph{Journal of Finance}, 44(1):205--209, 1989.
\bibitem{kalotay1993} A. Kalotay, G. Williams and F. Fabozzi, A model for valuing bonds and embedded options,
\emph{Financial Analysts Journal}, 49(3):35--46, 1993.
\bibitem{lando1998} D. Lando, On Cox processes and credit risky securities, \emph{Review of Derivatives Research},
2:99--120, 1998.
\bibitem{longstaff2001} F. Longstaff and E. Schwartz, Valuing American options by simulation: a simple least-squares
approach, \emph{Review of Financial Studies}, 14(1):113--147, 2001.
\bibitem{rannacher1984} R. Rannacher, Finite element solution of diffusion problems with irregular data,
\emph{Numerische Mathematik}, 43:309--327, 1984.
\end{thebibliography}

\end{document}
